% =====================================================================
% บทที่ 3 วิธีการดำเนินงาน (คู่มือข้อ 17.3)
% อธิบายขั้นตอนและวิธีการดำเนินงาน และเครื่องมือที่ใช้ในการทำโครงงาน
% *** ก่อนส่ง: ลบข้อความตัวอย่าง (...) ***
% =====================================================================

\section{ขั้นตอนการดำเนินงาน}
\label{sec:method}

โครงงานนี้ดำเนินงานตามแนวทาง\ldots{} (เช่น Agile แบบ Scrum หรือ Waterfall)
แบ่งเป็น 5 ขั้นตอน ดังภาพที่~\ref{fig:method}

% แผนภาพขั้นตอนวาดด้วย TikZ (แก้ข้อความในวงเล็บปีกกาได้โดยตรง)
% หรือใช้ภาพจากโปรแกรมอื่น: \includegraphics[width=0.9\linewidth]{images/ชื่อไฟล์}
\begin{figure}[H]
\centering
\begin{tikzpicture}[
    node distance=7bp and 10bp,
    stage/.style={draw=black!70, fill=black!4, rounded corners=3pt, align=center,
      font=\normalfont\small, minimum height=34bp, text width=64bp, inner sep=3bp},
    arr/.style={-{Stealth[length=5bp]}, line width=0.8pt, draw=black!60}]
  \node[stage] (s1) {ศึกษาปัญหาและ\\ความต้องการ};
  \node[stage, right=of s1] (s2) {วิเคราะห์และ\\ออกแบบระบบ};
  \node[stage, right=of s2] (s3) {พัฒนาระบบ};
  \node[stage, right=of s3] (s4) {ทดสอบและ\\ประเมินผล};
  \node[stage, right=of s4] (s5) {ปรับปรุงและ\\สรุปผล};
  \draw[arr] (s1) -- (s2); \draw[arr] (s2) -- (s3);
  \draw[arr] (s3) -- (s4); \draw[arr] (s4) -- (s5);
  \draw[arr] (s4.south) -- ++(0,-12bp) -| node[pos=0.25, below, font=\normalfont\footnotesize]{พบข้อผิดพลาด} (s3.south);
\end{tikzpicture}
\caption{ขั้นตอนการดำเนินงาน}
\label{fig:method}
\end{figure}

\subsection{ศึกษาปัญหาและรวบรวมความต้องการ}

สัมภาษณ์\ldots{} จำนวน\ldots{} คน และศึกษาเอกสาร\ldots{} เพื่อรวบรวมความต้องการของระบบ

\subsection{วิเคราะห์และออกแบบระบบ}

วิเคราะห์ความต้องการด้วยแผนภาพ UML และออกแบบฐานข้อมูล (รายละเอียดในบทที่ 4)

\subsection{พัฒนาระบบ}

พัฒนาระบบด้วยภาษา\ldots{} และเฟรมเวิร์ก\ldots{} โดยแบ่งการพัฒนาเป็นรอบละ\ldots{} สัปดาห์

\subsection{ทดสอบและประเมินผล}

ทดสอบระบบด้วยกรณีทดสอบ (Test case) และประเมินความพึงพอใจของผู้ใช้ด้วยแบบสอบถาม
โดยวิเคราะห์คะแนนด้วยค่าเฉลี่ยตามสมการที่~\eqref{eq:mean}

% สมการ: เลขสมการในวงเล็บชิดขวาให้อัตโนมัติ (1), (2), ... หรือ (3.1) เมื่อใช้ chapterfloats
%   \eqcaption{ชื่อสมการ} -> แสดงในสารบัญสมการ (สร้างหน้าสารบัญสมการให้เมื่อมีอย่างน้อย 1 สมการ)
%   อ้างถึงในเนื้อหาด้วย สมการที่~\eqref{ชื่อ label} -> "สมการที่ (1)"
%   ข้อความภาษาไทยในสมการเขียนใน \text{...}
\begin{equation}
\bar{x} = \frac{1}{n}\sum_{i=1}^{n} x_i
\eqcaption{ค่าเฉลี่ยเลขคณิต}
\label{eq:mean}
\end{equation}
% คำอธิบายสัญลักษณ์ใต้สมการ: คอลัมน์แรกเป็นสัญลักษณ์ (โหมดคณิตศาสตร์) คอลัมน์ที่สองเป็นคำอธิบาย
\begin{eqwhere}
\bar{x} & ค่าเฉลี่ยของคะแนน \\
x_i     & คะแนนของผู้ตอบคนที่ $i$ \\
n       & จำนวนผู้ตอบแบบสอบถาม
\end{eqwhere}

\subsection{ปรับปรุงและสรุปผล}

ปรับปรุงระบบตามผลการทดสอบ และจัดทำรายงานและคู่มือการใช้งาน

\section{เครื่องมือที่ใช้ในการดำเนินงาน}
\label{sec:tools}

\subsection{ฮาร์ดแวร์}

ฮาร์ดแวร์ที่ใช้ในการพัฒนาและทดสอบระบบแสดงดังตารางที่~\ref{tab:hardware}

\begin{table}[H]
\caption{ฮาร์ดแวร์ที่ใช้ในการพัฒนาและทดสอบระบบ}
\label{tab:hardware}
\centering
\begin{tabular}{|C{1.2cm}|L{4.4cm}|L{\dimexpr\linewidth-5.6cm-6\tabcolsep-4\arrayrulewidth\relax}|}
\hline
\textbf{ลำดับ} & \textbf{รายการ} & \textbf{รายละเอียด} \\ \hline
1 & คอมพิวเตอร์สำหรับพัฒนา & CPU \ldots{}, RAM \ldots{} GB, SSD \ldots{} GB \\ \hline
2 & เครื่องแม่ข่าย & \ldots{} \\ \hline
3 & โทรศัพท์มือถือสำหรับทดสอบ & \ldots{} \\ \hline
\end{tabular}
\end{table}

\subsection{ซอฟต์แวร์}

ซอฟต์แวร์ที่ใช้ในการพัฒนาระบบแสดงดังตารางที่~\ref{tab:software}

\begin{table}[H]
\caption{ซอฟต์แวร์ที่ใช้ในการพัฒนาระบบ}
\label{tab:software}
\centering
\begin{tabular}{|C{1.2cm}|L{4.2cm}|C{2.2cm}|L{\dimexpr\linewidth-7.6cm-8\tabcolsep-5\arrayrulewidth\relax}|}
\hline
\textbf{ลำดับ} & \textbf{ซอฟต์แวร์} & \textbf{เวอร์ชัน} & \textbf{การใช้งาน} \\ \hline
1 & Visual Studio Code & 1.9x & เขียนโปรแกรม \\ \hline
2 & Node.js & 20 LTS & สภาพแวดล้อมการทำงานของส่วนหลัง \\ \hline
3 & React & 18 & พัฒนาส่วนหน้า \\ \hline
4 & PostgreSQL & 15 & ระบบจัดการฐานข้อมูล \\ \hline
5 & Git และ GitHub & -- & จัดการเวอร์ชันของโค้ด \\ \hline
\end{tabular}
\end{table}

\section{แผนการดำเนินงาน}
\label{sec:plan}

แผนการดำเนินงานตลอดระยะเวลาโครงงานแสดงดังตารางที่~\ref{tab:gantt}

% แผนการดำเนินงาน (Gantt chart): ใส่ \cellcolor{ganttbar} ในช่องเดือนที่ทำกิจกรรม
\begin{table}[H]
\caption{แผนการดำเนินงาน}
\label{tab:gantt}
\centering
\begin{tabular}{|C{1.2cm}|L{\dimexpr\linewidth-7.8cm-16\tabcolsep-9\arrayrulewidth\relax}|*{6}{C{1.1cm}|}}
\hline
\textbf{ลำดับ} & \textbf{กิจกรรม} & \textbf{มิ.ย.} & \textbf{ก.ค.} & \textbf{ส.ค.}
  & \textbf{ก.ย.} & \textbf{ต.ค.} & \textbf{พ.ย.} \\ \hline
1 & ศึกษาปัญหาและรวบรวมความต้องการ & \cellcolor{ganttbar} & & & & & \\ \hline
2 & วิเคราะห์และออกแบบระบบ & & \cellcolor{ganttbar} & \cellcolor{ganttbar} & & & \\ \hline
3 & พัฒนาระบบ & & & \cellcolor{ganttbar} & \cellcolor{ganttbar} & \cellcolor{ganttbar} & \\ \hline
4 & ทดสอบและประเมินผล & & & & & \cellcolor{ganttbar} & \cellcolor{ganttbar} \\ \hline
5 & จัดทำรายงานและคู่มือ & & & & & & \cellcolor{ganttbar} \\ \hline
\end{tabular}
\end{table}
